% Source: https://artofproblemsolving.com/community/c3h1573693_area_of_a_triangle
% Stable identity: aops-topic-1573693-q01
The side from $(-\pi,\pi)$ to $(\pi,\pi)$ is horizontal (on the line $y=\pi$) and has length $2\pi$. The opposite vertex $(0,0)$ is at vertical distance $\pi$ from this line.

\begin{asy}
usepackage("amsmath");
size(5cm);
defaultpen(fontsize(10pt));
import graph;
real p=pi;
pair A=(0,0),B=(p,p),C=(-p,p);
draw((-4.2,0)--(4.2,0),gray,Arrow(TeXHead));draw((0,-0.4)--(0,4.9),gray,Arrow(TeXHead));
label("$x$",(4.2,0),E,gray);label("$y$",(0,4.9),N,gray);
fill(A--B--C--cycle,paleyellow);draw(A--B--C--cycle,linewidth(0.9));
dot(A);dot(B);dot(C);
label("$(0,0)$",A,SE);label("$(\pi,\pi)$",B,NE);label("$(-\pi,\pi)$",C,NW);
draw(A--(0,p),dashed);label("$\pi$",(0,p/2),W);draw((-p,p+0.55)--(p,p+0.55),gray,Arrows(TeXHead));label("$2\pi$",(-1.6,p+0.55),N);
label("\mbox{}",(min(currentpicture,user=true)+max(currentpicture,user=true))/2,black); // keeps the text after the figure black when LaTeX typesets the labels inline
\end{asy}

So the area is
\[
\dfrac{1}{2}\cdot2\pi\cdot\pi=\boxed{\pi^2}.
\]
